\addcontentsline{toc}{chapter}{Resumo}

\begin{center}
\huge{{\bf Resumo}}
\vspace{1cm}
\end{center}

\begin{otherlanguage}{portuguese}
MicroTEDAClus é um algoritmo de agrupamento evolutivo para fluxos de dados que resume o fluxo em microgrupos por meio da excentricidade e da tipicidade do arcabouço TEDA, e cujos autores relatam ter desempenho equivalente ao de algoritmos consolidados de agrupamento de fluxos e se adaptar melhor do que eles a mudanças nos dados. Este projeto final de curso, desenvolvido na Universidade Federal de Minas Gerais, no curso de Engenharia de Controle e Automação, sob orientação do Prof. Janier Arias Garcia, reavalia o algoritmo e propõe um modelo aprimorado. A reavaliação identifica problemas na formulação do algoritmo, diferenças entre a formulação e a implementação pública que produziu os seus resultados publicados e falhas na metodologia dos seus experimentos, entre elas etapas de pré-processamento que vazam informação sobre os dados e uma avaliação que prediz cada lote depois de aprendê-lo. Sob uma metodologia sem essas falhas, na qual todos os modelos são executados na mesma linguagem sobre os dados sem pré-processamento, avaliados de forma prequencial, ajustados pela mesma busca de hiperparâmetros e comparados por testes estatísticos em muitas execuções, o MicroTEDAClus não tem desempenho equivalente ao dos modelos de referência do trabalho original, nem se adapta melhor do que eles a mudanças de conceito.

O modelo aprimorado, MicroTEDAClus+, resulta de experimentos controlados nos quais cada correção ou modificação é aceita por uma regra definida previamente, em execuções separadas das da sua avaliação. Nos conjuntos de dados sintéticos do trabalho original, ele supera o MicroTEDAClus em qualidade de agrupamento, tempo de execução e memória, nenhum modelo de referência o supera nos conjuntos com mudança de conceito, e ele é o mais rápido desse experimento. Em um experimento modernizado no ambiente, nos modelos de referência e na seleção de hiperparâmetros, ele supera todos os modelos de referência em alguns conjuntos de dados e é superado por todos em outros. Em oito conjuntos de dados reais, porém, ele é superado pela maioria dos modelos de referência em seis. Um diagnóstico e novos experimentos controlados nesses conjuntos produzem a versão final do MicroTEDAClus+, que reduz a faixa de busca do seu limite de variância, abandona o rótulo de ruído e busca o fator que liga os seus microgrupos. Ela é superada pela maioria dos modelos de referência em três conjuntos de dados reais, em vez de seis, e supera todos em um, mas segue atrás de alguns nos fluxos reais com mudança de conceito.

\noindent\textbf{Palavras-chave:} Agrupamento de fluxos de dados. Agrupamento evolutivo. TEDA. Mudança de conceito. Metodologia experimental.
\end{otherlanguage}

\clearpage
\thispagestyle{empty}
\cleardoublepage
